\begin{center}
{\headingfont\LARGE\bfseries Preface} \\
\end{center}
\addcontentsline{toc}{section}{Preface}


The ability to design integrated circuits has become one of the defining technological capabilities of the modern world. From communication systems and transportation infrastructure to medical devices, defense electronics, artificial intelligence accelerators, and high-performance computing platforms, nearly every advanced engineering discipline today relies upon semiconductor technology. Yet the process of designing semiconductor devices is not enabled by design software alone. Behind every successful VLSI laboratory exists a carefully constructed computing infrastructure capable of supporting Electronic Design Automation (EDA) workflows reliably, consistently, and at scale. \\

The creation of such an environment is rarely visible to students. When a laboratory workstation is switched on and a design tool launches successfully, years of operating system development, network engineering, software integration, licensing infrastructure, and administrative planning remain hidden beneath the surface. However, establishing that environment from scratch is often a complex undertaking involving far more than the installation of a single software package. \\

This document originated from the practical challenges encountered during the establishment of the semiconductor design infrastructure used within the Department of Electronics and Communication Engineering at Gati Shakti Vishwavidyalaya. The work described throughout this manual was carried out as part of a broader effort to deploy a standardized Linux-based CAD environment capable of supporting Cadence design tools under the Chips-to-Startup (C2S) Programme. \\

The deployment process involved the installation and configuration of Red Hat Enterprise Linux workstations, preparation of user environments, integration of licensing services, workstation provisioning, software deployment, shell configuration, network validation, and the establishment of a repeatable administrative workflow. While many of these tasks appear straightforward when viewed individually, their interaction within a laboratory environment introduces a level of complexity that is often underappreciated until deployment begins. \\

A significant portion of this work was carried out by the author under the guidance of Dr. Prasiddh Trivedi, Assistant Professor, Department of Electrical Engineering, whose support, technical direction, and institutional coordination made the deployment possible. Throughout the setup process, numerous infrastructure decisions had to be evaluated and justified, ranging from operating system selection and workstation architecture to user naming conventions, deployment methodologies, storage allocation strategies, and license-compliance requirements. The resulting environment was not assembled through trial and error alone, but through a deliberate attempt to create a platform that could be maintained, reproduced, and expanded as the laboratory evolved. \\

One of the recurring observations during the deployment was the scarcity of documentation describing the complete setup process from an institutional perspective. Vendor manuals typically explain how to install a particular software package, while operating system documentation explains how to configure a particular service. What is often missing is a unified narrative describing how these components fit together within a real laboratory environment. Future administrators are therefore frequently forced to reconstruct deployment decisions by examining existing machines, consulting fragmented notes, or relying upon institutional memory. \\

This manual was written to address that problem. \\

Its purpose is not merely to record commands that were executed during installation, but to capture the reasoning behind those decisions. Wherever possible, procedures are accompanied by explanations describing why a particular configuration was chosen, what assumptions were made, and what consequences may arise if those assumptions change. The objective is to transform administrative knowledge from an informal collection of personal experience into a structured and reproducible reference. \\

The need for such documentation becomes especially important in academic environments. Students graduate, research projects conclude, hardware is replaced, operating systems reach end-of-life, and administrative responsibilities change hands. Without documentation, knowledge accumulated during deployment can disappear remarkably quickly. A workstation that functions perfectly today may become difficult to rebuild several years from now if the decisions that produced it were never recorded. By documenting both procedures and rationale, the laboratory reduces its dependence on individual administrators and preserves institutional knowledge for future use.\\

At the time of writing, many of the systems described in this document are already being deployed directly rather than reconstructed from documentation. Consequently, it is entirely possible that no immediate reader will ever need to follow every procedure presented here from beginning to end. In fact, the most successful outcome may be that future administrators rarely need this manual at all because the infrastructure continues to operate smoothly. Nevertheless, the value of technical documentation should not be measured solely by how often it is consulted. Its true purpose is to exist when something eventually changes: when a workstation fails, when a laboratory expands, when software must be reinstalled, or when a future administrator asks a simple but important question—"How was this system originally set up?"\\ 

This handbook therefore serves both as an operational reference and as a historical record of the deployment effort undertaken to establish the semiconductor CAD environment at Gati Shakti Vishwavidyalaya. It documents not only the final configuration of the systems, but also the engineering considerations, constraints, and administrative decisions that shaped their design.\\

The chapters that follow describe the deployment of the workstation environment, operating system preparation, configuration standards, and supporting infrastructure required for Cadence-based semiconductor design workflows. Together, they represent an attempt to ensure that the knowledge acquired during this deployment remains accessible, reproducible, and useful long after the initial installation work has been completed.\\ 
